\documentclass[12pt,a4paper]{article}
\usepackage[utf8]{inputenc}
\usepackage[T1]{fontenc}
\usepackage[brazilian]{babel}
\usepackage{mathptmx}
\usepackage[top=3cm,bottom=2.5cm,left=3cm,right=2.5cm]{geometry}
\usepackage{booktabs}
\usepackage{microtype}
\usepackage{url}
\usepackage[hidelinks]{hyperref}
\usepackage{enumitem}
\setlist{itemsep=2pt,topsep=4pt}
\usepackage{titlesec}
\titlespacing*{\section}{0pt}{12pt}{6pt}
\titlespacing*{\subsection}{0pt}{10pt}{4pt}

\title{\vspace{-1.2cm}\textbf{Assistente de banco de dados com GraphRAG: extração de
conhecimento, grafo no Neo4j, recuperação híbrida e respostas fundamentadas}}
\author{Autor 1, Autor 2, Autor 3, Autor 4, Autor 5, Autor 6\\
\small Instituição --- Curso --- Disciplina de Banco de Dados\\
\small \texttt{preencher com os nomes e o e-mail de cada integrante}}
\date{}

\begin{document}
\maketitle

\begin{abstract}
\noindent Assistentes que respondem perguntas sobre bancos de dados precisam de mais que
fluência: precisam de rastreabilidade. Este trabalho descreve um assistente que extrai conceitos
e relações de nove capítulos de material didático sobre bancos de dados relacionais e SQL,
armazena esse conhecimento em um grafo no Neo4j, recupera evidências por busca vetorial e por
travessia de grafo, e gera respostas em português com citação literal por afirmação e um nível de
sustentação calculado por regras, nunca pelo modelo. O sistema é avaliado em 15 perguntas: o modo
híbrido (vetor e grafo) acerta 93\% contra 60\% da busca puramente vetorial, diferença concentrada
nas perguntas que exigem dois ou mais relacionamentos (4 de 5 contra 1 de 5). Os três modos
abstiveram-se corretamente em todas as cinco perguntas fora do corpus, e nenhuma das 43 respostas
produzidas trouxe citação não verificada.

\vspace{4pt}
\noindent \textbf{Palavras-chave:} GraphRAG; KAG; Neo4j; recuperação híbrida; fundamentação de
respostas; LLM.
\end{abstract}

\section{Introdução}

Modelos de linguagem respondem com fluência sobre bancos de dados e SQL, mas fluência não é
rastreabilidade. Em contexto de ensino e de documentação técnica, uma resposta sem fonte
verificável tem o mesmo valor que uma resposta ausente: não se pode conferir, e não se pode
ensinar a conferir.

O projeto parte de uma pergunta prática --- \emph{a pergunta feita na interface consegue ser
respondida com informação de dois documentos reais, com fonte verificável e visualização do
grafo?} --- e a trata como restrição de projeto, não como aspiração. O sistema foi construído em
seis etapas encadeadas: extração de conhecimento do corpus, modelagem e ingestão em um grafo,
recuperação híbrida, geração validada por LLM, interface com visualização e avaliação comparativa.

As contribuições são: (i) um contrato de extração com vocabulário fechado de rótulos, tipos de
relação e propriedades, que permite validar automaticamente tudo o que entra no grafo; (ii) uma
recuperação híbrida que soma busca vetorial multilíngue e travessia de grafo por fusão de ranks,
com caminho mais curto e centralidade de grau; (iii) uma camada de geração em que o modelo
\emph{não} escreve Cypher: ele escolhe um plano em vocabulário fechado, e o sistema compila e
valida a consulta; (iv) uma verificação de resposta que confere cada citação contra o texto
original do corpus e calcula o nível de sustentação por regras explícitas; e (v) uma avaliação
comparativa com 15 perguntas, medida nos três modos de recuperação.

\section{Fundamentação e trabalhos relacionados}

\textbf{RAG.} A geração aumentada por recuperação \cite{lewis2020} combina um índice de passagens
com um modelo gerador, reduzindo a dependência de conhecimento paramétrico. A limitação que motiva
este trabalho é conhecida: a similaridade textual recupera passagens \emph{parecidas} com a
pergunta, não a cadeia de fatos que a responde.

\textbf{GraphRAG e KAG.} Abordagens que organizam o corpus em grafo \cite{edge2024} mostram ganho
em perguntas que exigem agregar informação dispersa. A ideia central adotada aqui é a mesma ---
conhecimento estruturado como intermediário --- com uma diferença de escopo: empregamos um grafo
\emph{curado e com evidência por aresta}, em vez de um grafo construído inteiramente pelo modelo,
o que permite verificar cada afirmação contra o texto de origem.

\textbf{Fundamentação e abstenção.} A literatura sobre alucinação \cite{ji2023} distingue erros de
fato de erros de formato e recomenda verificação externa. Nossa verificação é mecânica: a citação
precisa existir literalmente no trecho, e a resposta que cita identificador inexistente é
recusada.

\textbf{Fusão de rankings.} O \emph{Reciprocal Rank Fusion} \cite{cormack2009} combina rankings de
recuperadores distintos sem calibrar escores, e é usado aqui para unir busca vetorial e evidências
de grafo. O modelo de dados segue o grafo de propriedades \cite{robinson2015}, com três rótulos,
relações tipadas, índices de texto e um índice vetorial nativo do Neo4j 5.

\section{Corpus e extração de conhecimento}

\subsection{Corpus}

Nove documentos em inglês, de material didático aberto de banco de dados relacional e da
documentação do PostgreSQL, com autores, licença e URL preservados por documento: D01 (conceitos
fundamentais), D02 (modelo ER), D03 (integridade e restrições), D04 (dependências funcionais), D05
(normalização), D06 (consulta e filtro), D07 (junções), D08 (funções de agregação) e D09 (modelo
relacional).

O corpus foi dividido em \textbf{59 trechos}: 31 trechos de texto de página e 28 suplementos que
transcrevem figuras essenciais, marcados com estado de transcrição e incertezas declaradas, para
que nenhuma transcrição seja confundida com texto original. Cada trecho carrega documento, caminho
de seção e URL de origem.

\subsection{Extração}

Três métodos, declarados por conceito e por relação:

\begin{table}[h]
\centering
\small
\begin{tabular}{@{}llrr@{}}
\toprule
Arquivo & Método & Conceitos & Relações \\
\midrule
\texttt{rule-knowledge-union-v2} & regras sobre o texto, com curadoria & 79 & 5 \\
\texttt{model-knowledge-union-v3} & proposta de LLM local, com curadoria & 20 & 11 \\
\texttt{manual-source-review-v1} & curadoria manual assistida por IA & 5 & 0 \\
\bottomrule
\end{tabular}
\end{table}

Os três arquivos somam \textbf{93 conceitos únicos} e \textbf{16 triplas semânticas}; com a amostra
manual opcional, \textbf{23 triplas}. Cada conceito tem nome canônico em inglês, nome em português,
aliases, descrição, estado de revisão e as citações que o sustentam. Cada relação carrega os
identificadores dos trechos de evidência e as citações literais correspondentes: a unidade de
verificação de todo o sistema é a aresta com evidência, não o documento.

Identificadores são estáveis e independentes do nome (\texttt{C\_\ldots} para conceitos,
\texttt{D03\_C001} para trechos, \texttt{D0n} para documentos). Nome nunca é identidade, o que
evita fundir conceitos por semelhança de rótulo.

\subsection{Limites desta etapa}

A cobertura de relações é o gargalo do projeto. Vinte e três triplas semânticas para 93 conceitos
significam que a maioria dos conceitos não participa de nenhum caminho multi-hop. Dos três
métodos, apenas o de regras e o de LLM produzem relações, e nenhum foi validado por revisão humana
independente: o registro é de aprovação do próprio grupo.

\section{Modelagem e ingestão no Neo4j}

\subsection{Modelo}

Três rótulos, cada um com identificador único e estável, e relações estruturais separadas das
semânticas. Estruturais: \texttt{CONTAINS} (Documento para Trecho), \texttt{SUPPLEMENTS} (figura
para trecho-pai) e \texttt{MENTIONS} (Trecho para Conceito). Semânticas, em vocabulário fechado de
14 tipos: \texttt{IS\_A}, \texttt{SUPPORTS}, \texttt{REQUIRES}, \texttt{GOVERNS},
\texttt{TYPICALLY\_USED\_IN}, \texttt{COMBINES}, \texttt{HAS\_FORM}, \texttt{PRESERVES},
\texttt{RETRIEVES}, \texttt{FILTERS}, \texttt{GROUPS}, \texttt{AGGREGATES},
\texttt{EVALUATED\_BEFORE} e \texttt{REMOVES\_DUPLICATES}.

O Neo4j armazena apenas valores simples, então objetos e listas aninhadas (evidências, licença,
caminhos de seção) ficam como texto JSON em propriedades terminadas em \texttt{\_json}, o que
mantém a ingestão idempotente e preserva a citação literal.

\subsection{Ingestão e garantias}

A ingestão roda em uma única transação, com \texttt{MERGE} por identificador e \texttt{SET +=}
para propriedades, de modo que uma segunda execução não duplica nem apaga --- inclusive
propriedades criadas por outras etapas, como o vetor de \emph{embedding}. Três restrições de
unicidade e sete índices, entre eles um índice de texto sobre nomes em inglês, português e
aliases que ignora acentos.

\begin{table}[h]
\centering
\small
\begin{tabular}{@{}lrrrrr@{}}
\toprule
Carga & Document & Chunk & Concept & MENTIONS & Rel. semânticas \\
\midrule
3 arquivos principais & 9 & 59 & 93 & 130 & 16 \\
com a amostra opcional & 9 & 59 & 93 & 134 & 23 \\
\bottomrule
\end{tabular}
\end{table}

Uma verificação automática confere o banco contra os arquivos entregues: contagens, ausência de
referência inválida, presença de caminho de dois saltos com evidência por aresta e ausência de
\texttt{REQUIRES} onde a fonte diz ``tipicamente''.

\section{Recuperação híbrida}

\subsection{Embeddings}

Modelo \texttt{paraphrase-multilingual-MiniLM-L12-v2}, 384 dimensões, similaridade de cosseno. A
escolha responde ao descompasso de idioma do projeto: o corpus é em inglês e as perguntas em
português. Os 59 vetores estão pré-computados e versionados, e o mesmo modelo cria um índice
vetorial nativo no Neo4j.

\subsection{Componentes}

\begin{enumerate}
\item \textbf{Vinculação de conceitos.} A pergunta é comparada com os 93 conceitos (nome em
  inglês, em português e aliases) após normalização sem acento e sem plural.
\item \textbf{Travessia multi-hop.} A partir das sementes, o sistema percorre as relações
  semânticas até um orçamento de saltos, carregando de cada aresta a citação que a sustenta.
\item \textbf{Caminho mais curto.} Quando a pergunta envolve dois conceitos, o caminho mínimo
  entre eles define a cadeia a percorrer --- é o que liga documentos distintos, como D03 e D07 na
  cadeia integridade referencial, chave estrangeira, junção, junção externa, junção externa à
  esquerda.
\item \textbf{Centralidade de grau.} Pondera os conceitos do subgrafo para priorizar nós de maior
  conectividade contextual.
\item \textbf{Fusão por RRF.} Os rankings vetorial e de grafo são combinados por uma soma de
  recíprocos das posições, com constante 60, o que soma os dois canais sem calibrar escores.
\end{enumerate}

\subsection{Modos de recuperação}

\begin{table}[h]
\centering
\small
\begin{tabular}{@{}lll@{}}
\toprule
Modo & Recuperação & Depende de \\
\midrule
\texttt{vector} & busca vetorial sobre os 59 trechos (\emph{top\_k} = 5) & embeddings \\
\texttt{hybrid} & vetorial, travessia, caminho mais curto e centralidade, por RRF & idem \\
\texttt{graph} & plano de consulta Cypher validado, executado no Neo4j & grafo ingerido \\
\bottomrule
\end{tabular}
\end{table}

\section{Geração validada com LLM}

\subsection{O modelo não escreve a consulta}

O modelo recebe a pergunta e devolve um \emph{plano} em vocabulário fechado: nó inicial, direção,
tipos de relação e limite. O sistema compila esse plano em uma única consulta de caminho,
parametrizada. Toda identificação --- rótulo, tipo, propriedade, identificador --- é conferida
contra listas fixas antes de chegar ao Cypher, de modo que texto do modelo nunca vira sintaxe.

O orçamento de saltos é um único número (\texttt{DEFAULT\_MAX\_HOPS = 4}) compartilhado pelo
validador, pelo executor e pelos adaptadores. Cypher livre existe como alternativa, validado por
regras: uma única instrução, nenhum comentário, nenhuma string literal (valores só por parâmetro),
nenhum comando de escrita ou procedimento de extensão, todo padrão de nó rotulado,
\texttt{RETURN} projetado e \texttt{LIMIT} obrigatório. Um conjunto de 38 casos de contorno (32
recusas obrigatórias, 5 controles e um limite conhecido) é executado a cada rodada de testes.

\subsection{Resposta fundamentada}

O modelo recebe as evidências com identificadores (E1 \ldots\ En), nunca URLs, e devolve
afirmações, cada uma citando identificadores. O sistema então recusa citação de identificador
inexistente, confere se a citação existe \emph{literalmente} no texto do trecho, descarta
afirmações sem evidência, resolve cada citação em trecho, seção e URL, e calcula o nível de
sustentação por regras.

\begin{table}[h]
\centering
\small
\begin{tabular}{@{}ll@{}}
\toprule
Nível & Critério \\
\midrule
\texttt{high} & citações aceitas, de ao menos dois trechos e dois documentos \\
\texttt{medium} & um único trecho, um único documento, ou citação não conferida \\
\texttt{low} & evidência pendente ou sinalizada, ou afirmação descartada \\
\texttt{insufficient} & nenhuma afirmação sobreviveu à verificação \\
\bottomrule
\end{tabular}
\end{table}

O modelo é proibido de informar confiança: o nível é calculado fora dele, sobre as citações
efetivamente devolvidas.

\section{Interface e visualização}

A aplicação apresenta, na mesma tela: a resposta em português; cada afirmação com o marcador da
evidência que a sustenta; o painel de sustentação (nível, cobertura, trechos, documentos, citações
não verificadas e as razões); o subgrafo percorrido; as fontes com link; as citações literais; e a
consulta Cypher executada com seus parâmetros.

O subgrafo é desenhado em SVG com layout por distância na travessia, determinístico: a mesma
resposta sempre produz a mesma figura. No modo híbrido, o grafo é montado a partir das relações que
o pacote de evidências realmente usou, e a tela informa canais e origens em vez de exibir um
Cypher que não foi executado.

Como a geração local leva de um a dois minutos por resposta, a execução é assíncrona: a requisição
devolve um identificador de trabalho e a página acompanha o passo corrente. Estados terminais são
distintos de propósito --- respondido, sem correspondência no grafo, resposta recusada pela
verificação e falha de serviço --- e nenhum deles é suavizado em texto.

\section{Avaliação}

\subsection{Desenho}

Quinze perguntas em três grupos de cinco: conceituais, multi-hop e sem resposta no corpus. Cada
pergunta tem gabarito de conceitos e de documentos, \emph{derivado das evidências das relações do
grafo} --- nunca das respostas do sistema. As perguntas multi-hop foram escolhidas entre caminhos
que existem, e um teste automatizado confere que cada par consecutivo da cadeia tem aresta de um
salto.

O critério de acerto é mecânico e idêntico nos três modos: pergunta com resposta exige resposta
entregue, ao menos um documento esperado entre as fontes citadas e nenhuma citação não verificada;
pergunta sem resposta exige abstenção. O mesmo modelo, prompt, orçamento de tokens e verificador
são usados; muda apenas a recuperação.

\subsection{Resultados}

Quarenta e cinco execuções (15 perguntas em 3 modos), 2 falhas de verificação e nenhuma citação
não verificada.

\begin{table}[h]
\centering
\small
\begin{tabular}{@{}lrrrrrr@{}}
\toprule
Modo & Corretas & Acerto & Doc.\ esperado & Abstenção & Nível alto & Tempo \\
\midrule
\texttt{vector} & 9/15 & 60\% & 4/10 & 5/5 & 0 & 65 s \\
\texttt{hybrid} & 14/15 & \textbf{93\%} & 9/10 & 5/5 & 5 & 71 s \\
\texttt{graph} & 14/15 & \textbf{93\%} & 9/10 & 5/5 & 1 & 51 s \\
\bottomrule
\end{tabular}
\end{table}

\begin{table}[h]
\centering
\small
\begin{tabular}{@{}lrrr@{}}
\toprule
Grupo & \texttt{vector} & \texttt{hybrid} & \texttt{graph} \\
\midrule
conceitual (5) & 3/5 & 5/5 & 4/5 \\
multi-hop (5) & 1/5 & 4/5 & 5/5 \\
sem resposta (5) & 5/5 & 5/5 & 5/5 \\
\bottomrule
\end{tabular}
\end{table}

\subsection{Análise}

\textbf{A diferença aparece exatamente onde a arquitetura promete.} No grupo multi-hop o baseline
vetorial acerta 1 de 5; o híbrido, 4 de 5; o modo de plano, 5 de 5. Em três dessas perguntas o
baseline recebeu trechos sobre o assunto e mesmo assim absteve-se: um trecho isolado não responde
a uma pergunta que exige ligar dois conceitos. É precisamente o que a travessia do grafo
acrescenta, e é a evidência que justifica a complexidade da arquitetura.

\textbf{Abstenção é 5 de 5 nos três modos.} Nenhum modo respondeu pergunta fora do corpus. No modo
de plano a recusa acontece antes do modelo: se nenhum conceito da pergunta existe no grafo, o
modelo não é chamado, o que também explica o menor tempo médio.

\textbf{O modo de plano iguala o acerto e troca diversidade por velocidade.} Ele produz uma única
resposta \texttt{high}, contra cinco do híbrido, porque um caminho tende a citar um documento só.

\textbf{Duas falhas de verificação, do mesmo tipo.} Em duas células o modelo devolveu abstenção
acompanhada de afirmações com evidência, e o verificador recusou. A resposta foi barrada em vez de
entregue inconsistente; para o critério, conta como não acerto.

\section{Ameaças à validade}

\textbf{Cobertura da extração.} Vinte e três relações semânticas condicionam tudo: a maior parte
das perguntas multi-hop possíveis não tem caminho no grafo, e as perguntas da avaliação usam
praticamente todos os caminhos disponíveis.

\textbf{Circularidade parcial do gabarito.} Os documentos esperados vêm das evidências do grafo, o
que favorece os modos que usam o grafo por construção. O critério também não julga a correção do
\emph{texto} da resposta, apenas a fonte citada e a literalidade da citação.

\textbf{Tamanho da amostra.} Quinze perguntas significam 6,7 pontos percentuais por pergunta, e
cada célula foi executada uma única vez. A mesma pergunta, com a mesma consulta, já produziu
níveis diferentes em execuções repetidas, porque o nível depende de quais evidências o modelo cita.

\textbf{Uma correção de gabarito declarada.} A pergunta Q09 passou a aceitar o documento D08 além
do D06 depois da primeira rodada, porque D08 também explica \texttt{WHERE}. Sem a correção,
\texttt{vector} teria 8/15 e os outros dois 13/15: a conclusão não muda.

\textbf{Revisão humana.} Nenhuma etapa teve revisão humana independente dos fatos extraídos, o que
enfraquece a premissa de que as arestas do grafo estão corretas.

\textbf{Um modelo, uma máquina.} \texttt{qwen3:8b} em CPU parcial; outro modelo mudaria os
números. A verificação das citações, não.

\section{Conclusão e continuidade}

O sistema cumpre o objetivo inicial: uma pergunta em português é respondida com informação de
documentos distintos, com citação literal por afirmação, nível de sustentação calculado por regras
e visualização do subgrafo percorrido. A comparação com RAG puramente vetorial mostra ganho
concentrado nas perguntas que exigem mais de um relacionamento --- 4 de 5 contra 1 de 5 --- e
nenhum modo respondeu pergunta fora do corpus.

Três continuidades se destacam: ampliar a extração de relações, que é o fator limitante de
cobertura; repetir a avaliação três vezes por célula e reportar média e dispersão; e substituir o
gabarito derivado do grafo por um gabarito anotado por pessoa, o que permitiria julgar a correção
do texto e não apenas a procedência da fonte.

\begin{thebibliography}{9}
\bibitem{lewis2020} LEWIS, P. et al. Retrieval-augmented generation for knowledge-intensive NLP
tasks. In: \emph{Advances in Neural Information Processing Systems (NeurIPS)}, 2020.
\bibitem{edge2024} EDGE, D. et al. From local to global: a graph RAG approach to query-focused
summarization. \emph{arXiv preprint arXiv:2404.16130}, 2024.
\bibitem{cormack2009} CORMACK, G. V.; CLARKE, C. L. A.; BUETTCHER, S. Reciprocal rank fusion
outperforms Condorcet and individual rank learning methods. In: \emph{Proceedings of SIGIR}, 2009.
p. 758--759.
\bibitem{ji2023} JI, Z. et al. Survey of hallucination in natural language generation.
\emph{ACM Computing Surveys}, v. 55, n. 12, p. 1--38, 2023.
\bibitem{robinson2015} ROBINSON, I.; WEBBER, J.; EIFREM, E. \emph{Graph Databases}. 2. ed.
Sebastopol: O'Reilly Media, 2015.
\end{thebibliography}

\end{document}
